\documentclass[11pt,a4paper]{article}
\usepackage[utf8]{inputenc}
\usepackage[T1]{fontenc}
\usepackage[margin=2cm, top=2.4cm]{geometry}
\usepackage{xcolor}
\usepackage{enumitem}
\usepackage{titlesec}
\usepackage{fancyhdr}
\usepackage{hyperref}
\usepackage{booktabs}
\usepackage{tabularx}
\usepackage{verbatim}
\usepackage{amsmath,amssymb}

% Course color palette
\definecolor{mlblue}{HTML}{0066CC}
\definecolor{mlpurple}{HTML}{3333B2}
\definecolor{mllavender}{HTML}{ADADE0}
\definecolor{mllavender2}{HTML}{C1C1E8}
\definecolor{mllavender3}{HTML}{CCCCEB}
\definecolor{mllavender4}{HTML}{D6D6EF}
\definecolor{mlorange}{HTML}{FF7F0E}
\definecolor{mlgreen}{HTML}{2CA02C}
\definecolor{mlred}{HTML}{D62728}
\definecolor{mlgray}{HTML}{7F7F7F}
\definecolor{lightgray}{HTML}{F0F0F0}
\definecolor{midgray}{HTML}{B4B4B4}

% Section heading colors
\titleformat{\section}{\large\bfseries\color{mlpurple}}{}{0em}{}[\vspace{-4pt}\textcolor{mllavender}{\rule{\linewidth}{1pt}}\vspace{2pt}]
\titleformat{\subsection}{\normalsize\bfseries\color{mlblue}}{}{0em}{\textbullet~}
\titlespacing*{\section}{0pt}{8pt}{4pt}
\titlespacing*{\subsection}{0pt}{5pt}{2pt}

% Header/footer
\pagestyle{fancy}
\fancyhf{}
\setlength{\headheight}{22pt}
\addtolength{\topmargin}{-10pt}
\renewcommand{\headrulewidth}{0.4pt}
\renewcommand{\headrule}{\hbox to\headwidth{\color{mllavender}\leaders\hrule height \headrulewidth\hfill}}
\fancyhead[L]{\small\color{mlpurple}\textbf{Introduction to Smart Contracts} \textcolor{mlgray}{|} Lesson 14 \textcolor{mlgray}{|} Pre-Class Discovery Handout}
\fancyhead[R]{\small\color{mlgray}Prof.\ Dr.\ Joerg Osterrieder}
\fancyfoot[C]{\small\color{mlgray}\thepage}

% Compact list spacing
\setlist[itemize]{nosep, topsep=2pt, partopsep=0pt, leftmargin=1.4em}
\setlist[enumerate]{nosep, topsep=2pt, partopsep=0pt, leftmargin=1.6em}

% Activity box helper
\newcommand{\activitybox}[3]{%
  \noindent\colorbox{mllavender4}{\parbox{\dimexpr\linewidth-2\fboxsep}{%
    \textbf{\color{mlpurple}#1}\hfill\textcolor{mlgray}{\small #2}\\[2pt]#3%
  }}\par\vspace{4pt}%
}

% Thin ruled cell for fill-in tables
\newcommand{\fillcell}{\rule{2.2cm}{0.3pt}}

\begin{document}

\begin{center}
{\LARGE\color{mlpurple}\textbf{Introduction to Smart Contracts}}\\[4pt]
{\large\color{mlblue}Pre-Class Discovery Handout}\\[2pt]
{\small\color{mlgray}Lesson 14 $\cdot$ Complete before class $\cdot$ 25--30 minutes}
\end{center}
\vspace{4pt}

% -- Activity 1 ----------------------------------------------------------------
\activitybox{Activity 1: Explore a Smart Contract on Etherscan}{10 min $\cdot$ Individual}{%
Visit \href{https://etherscan.io}{etherscan.io} and navigate to a well-known smart contract such as the USDC token contract (address: \texttt{0xA0b86991c6218b36c1d19D4a2e9Eb0cE3606eB48}). Explore the contract page and answer:
\begin{enumerate}
\item How many transactions has this contract processed? What is the total number of token holders?
\item Click the ``Contract'' tab and then ``Read Contract.'' Find the \texttt{name()}, \texttt{symbol()}, and \texttt{decimals()} functions. What values do they return?
\item Click ``Write Contract.'' What functions are available? Why would you need to connect a wallet to use them?
\item Look at the ``Events'' tab. What event is emitted most frequently? What information does a \texttt{Transfer} event contain (from, to, value)?
\end{enumerate}
\textbf{Bonus:} Find another smart contract (e.g., Uniswap Router) and compare its complexity (number of functions, transaction volume) to the USDC contract.
}

\vspace{2pt}

% -- Activity 2 ----------------------------------------------------------------
\activitybox{Activity 2: The Vending Machine Analogy}{5 min $\cdot$ Pairs}{%
Nick Szabo described a vending machine as the simplest form of a smart contract: you insert money, the machine checks the amount, and it dispenses the product---no negotiation, no trust, no intermediary.
\begin{enumerate}
\item List three real-world processes (besides vending machines) that work like a smart contract: deterministic input, automatic verification, guaranteed output.
\item For each process, identify: (a) what is the ``input,'' (b) what condition is checked, and (c) what is the automatic output.
\item Which of your examples could actually be implemented as a blockchain smart contract? What would be gained or lost?
\end{enumerate}
\vspace{4pt}
{\footnotesize\renewcommand{\arraystretch}{1.35}%
\setlength{\tabcolsep}{5pt}%
\begin{tabular}{@{}p{3.0cm}p{3.0cm}p{3.0cm}p{3.5cm}@{}}
\toprule
\textbf{Process} & \textbf{Input} & \textbf{Condition Checked} & \textbf{Automatic Output} \\
\midrule
\fillcell & \fillcell & \fillcell & \fillcell \\
\fillcell & \fillcell & \fillcell & \fillcell \\
\fillcell & \fillcell & \fillcell & \fillcell \\
\bottomrule
\end{tabular}}
\vspace{2pt}
}

\vspace{2pt}

% -- Activity 3 ----------------------------------------------------------------
\activitybox{Activity 3: Compare Smart Contract Platforms}{10 min $\cdot$ Small Groups}{%
Research and compare three smart contract platforms: Ethereum, Solana, and one platform of your choice (e.g., Cardano, Polkadot, Avalanche, or Arbitrum). Answer:
\begin{enumerate}
\item What programming language does each platform use for smart contracts?
\item What is the approximate transaction speed (TPS) and average transaction cost for each?
\item What consensus mechanism does each use? How does this affect security and decentralization?
\item Which platform would you choose to deploy a high-frequency trading DApp? A community governance DAO? Justify each choice.
\end{enumerate}
\vspace{4pt}
{\footnotesize\renewcommand{\arraystretch}{1.35}%
\setlength{\tabcolsep}{5pt}%
\begin{tabular}{@{}p{2.4cm}p{2.4cm}p{2.4cm}p{2.0cm}p{2.8cm}@{}}
\toprule
\textbf{Platform} & \textbf{Language} & \textbf{TPS / Cost} & \textbf{Consensus} & \textbf{Best For} \\
\midrule
Ethereum         & \fillcell & \fillcell & \fillcell & \fillcell \\
Solana           & \fillcell & \fillcell & \fillcell & \fillcell \\
\fillcell        & \fillcell & \fillcell & \fillcell & \fillcell \\
\bottomrule
\end{tabular}}
\vspace{2pt}
}

\vspace{2pt}

% -- Activity 4 ----------------------------------------------------------------
\activitybox{Activity 4: Smart Contract Use Case Design}{5 min $\cdot$ Individual}{%
Design a simple smart contract for one of the following scenarios (or invent your own):
\begin{itemize}
\item A crowdfunding contract that refunds donors if a goal is not met by a deadline
\item An escrow contract for freelance work that releases payment upon delivery confirmation
\item A lottery contract where participants buy tickets and a winner is selected
\end{itemize}
\begin{enumerate}
\item What \textbf{state variables} would your contract need? (e.g., owner address, balance, deadline)
\item What \textbf{functions} would users call? (e.g., \texttt{contribute()}, \texttt{withdraw()}, \texttt{refund()})
\item What \textbf{conditions} must be checked before executing each function? (e.g., ``deadline has not passed,'' ``caller is the owner'')
\item What could go wrong? Identify at least one security risk and how you would mitigate it.
\end{enumerate}
\vspace{4pt}
\noindent
\textbf{Contract name:}~\fillcell\quad
\textbf{Number of functions:}~\fillcell\quad
\textbf{Main security risk:}~\fillcell
\vspace{2pt}
}

\vspace{4pt}

% -- Glossary ------------------------------------------------------------------
\section{Key Terms}

\renewcommand{\arraystretch}{1.25}
\noindent\begin{tabular}{@{}p{3.6cm}p{11.0cm}@{}}
\toprule
\textbf{Term} & \textbf{Definition} \\
\midrule
\textbf{Smart Contract}  & A self-executing program stored on a blockchain that automatically enforces the terms of an agreement when predetermined conditions are met, without requiring intermediaries. \\
\textbf{Solidity}         & The most widely used programming language for writing smart contracts on Ethereum and EVM-compatible blockchains. Statically typed, influenced by JavaScript and C++. \\
\textbf{EVM}              & Ethereum Virtual Machine. The runtime environment that executes smart contract bytecode on every node in the Ethereum network, ensuring deterministic computation. \\
\textbf{Gas}              & A unit measuring the computational effort required to execute operations on Ethereum. Users pay gas fees (in ETH) to compensate validators for processing transactions. \\
\textbf{Gwei}             & A denomination of ETH equal to $10^{-9}$ ETH (one billionth). Gas prices are typically quoted in Gwei (e.g., 20 Gwei per unit of gas). \\
\textbf{Transaction}      & A signed message sent from an externally owned account to the blockchain. Transactions can transfer ETH, deploy contracts, or call contract functions. \\
\textbf{Block}            & A batch of validated transactions bundled together and appended to the blockchain. Each block references the previous block's hash, forming an immutable chain. \\
\textbf{Wallet}           & Software (e.g., MetaMask) or hardware (e.g., Ledger) that manages a user's private keys and enables signing transactions to interact with the blockchain. \\
\textbf{Private Key}      & A secret 256-bit number that proves ownership of a blockchain address and authorizes transactions. Must never be shared; loss means permanent loss of funds. \\
\textbf{Address}          & A 20-byte (160-bit) identifier derived from a public key, used to send and receive ETH and interact with smart contracts. Displayed as a 42-character hexadecimal string (0x\ldots). \\
\textbf{Deploy}           & The process of sending a transaction containing compiled smart contract bytecode to the blockchain, creating a new contract instance at a unique address. \\
\textbf{ABI}              & Application Binary Interface. A JSON specification describing a smart contract's functions, inputs, outputs, and events, enabling external applications to interact with it. \\
\textbf{DApp}             & Decentralized Application. A front-end application (web or mobile) that interacts with one or more smart contracts on a blockchain instead of a centralized server. \\
\textbf{Oracle}           & A service (e.g., Chainlink) that feeds external, off-chain data (prices, weather, sports scores) into smart contracts, which cannot natively access data outside the blockchain. \\
\textbf{Immutable}        & The property that once a smart contract is deployed to the blockchain, its code cannot be modified or deleted, ensuring permanence but requiring careful pre-deployment testing. \\
\bottomrule
\end{tabular}

\vspace{6pt}
\noindent\textcolor{mlgray}{\small\textit{Prepared by Prof.\ Dr.\ Joerg Osterrieder \,\textbullet\, Introduction to Smart Contracts --- Lesson 14 \,\textbullet\, Pre-Class Discovery Handout}}

\end{document}
